\chapter{Anomaly as surprise: the Gaussian, Mahalanobis distance and the RX detector}\label{ch:rx}

\begin{onething}
Every detector in this book answers one question, ``how badly does this pixel fit the pattern of
normal?'', and the RX detector answers it by modelling normal as one Gaussian cloud and scoring a
pixel by its Mahalanobis distance from the centre: Euclidean distance after whitening. On 165
correlated bands the covariance that whitening divides by is nearly singular, the distances explode,
and that is why plain RX was removed from the product on 2026-05-11.
\end{onething}

\section{Anomaly as surprise}\label{sec:rx:surprise}

\Cref{ch:problem} turned ``is this a gas leak?'' into ``is this pixel unlike its surroundings?''.
This chapter makes that computable. Every detector in this repository, statistical or neural,
follows the same three steps:
\begin{enumerate}
  \item build a model of what \emph{normal} looks like;
  \item measure how badly each pixel \emph{fits} that model;
  \item rank the pixels; the worst-fitting are the candidates.
\end{enumerate}
The two families of detectors differ only in step 1. The \emph{statistical} family describes normal
with statistics computed from the scene itself (its mean spectrum, and how its bands vary
together) and needs no training, no checkpoint and no labels. The \emph{reconstruction} family
(\cref{ch:recon}) trains a network to redraw ordinary imagery and treats a bad redrawing as a sign
of something unusual. This chapter and the next build the statistical family, and with it every
piece of mathematics the residual score of \cref{ch:residualmaha} needs.

Every detector outputs a \emph{score map}\index{score map}: one number per pixel, the height and
width of the scene, higher meaning more anomalous. It does not output a yes/no decision. A threshold
turns scores into a mask later, chosen by a person in a separate Action (\cref{ch:product}), and
invalid pixels are excluded from scoring and from every statistic computed over the map.

\subsection{Start in one dimension}

A thermal scene from Landsat~9 has one number per pixel, a temperature. The shipped thermal
statistics are a mean of 24.58\,\textdegree C\tagM{app/constants/thermal_pixel_consts.json} and a
standard deviation of 13.57\,\textdegree C\tagM{app/constants/thermal_pixel_consts.json}. How unusual
is a pixel at 71\,\textdegree C? Count how many standard deviations it sits above the mean:
\begin{align*}
z &= \frac{71 - 24.58}{13.57}\\
  &= \frac{46.42}{13.57}\\
  &= 3.42 .
\end{align*}
Three and a half standard deviations above the mean is unusual. That is the whole idea, and the
rest of the chapter generalises it to many bands. To generalise it properly you need to know
\emph{how} unusual 3.42 is, which needs a model of how the values are spread: the Gaussian.

\begin{analogy}
A doctor reading a blood panel does not flag a test because its number is large; she flags it
because it is large \emph{for that test}, measured against the normal range. When several tests move
together for an ordinary reason (dehydration raises many concentrations at once), she reads them as
one finding, not many. One value far outside its own normal range while the rest are ordinary is the
alarm. The RX detector reads a spectrum the same way: it measures each pixel against the normal
range of the scene, it discounts movements that the bands always make together, and it raises the
alarm for a movement the scene does not normally make.
\end{analogy}

\section{The one-dimensional Gaussian and the z-score}\label{sec:rx:gauss1d}

\begin{definition}[Gaussian distribution]\index{Gaussian distribution}
A quantity $x$ has a \emph{Gaussian} (or \emph{normal}) distribution with mean $\mu$ and standard
deviation $\sigma > 0$, written $x \sim \mathcal{N}(\mu, \sigma^2)$, if its probability density is
\begin{equation}\label{eq:rx:gauss1d}
p(x) = \frac{1}{\sqrt{2\pi}\,\sigma}\exp\!\left(-\frac{(x - \mu)^2}{2\sigma^2}\right).
\end{equation}
\end{definition}

A \emph{probability density}\index{probability density} is a curve whose area over an interval is
the probability of landing in that interval; the total area under it is 1. The Gaussian is the
familiar bell: highest at $\mu$, symmetric, falling away fast. Two pieces of \cref{eq:rx:gauss1d}
deserve a sentence each. The factor $1/(\sqrt{2\pi}\sigma)$ is only there to make the total area 1;
it does not depend on $x$. All the shape is in the exponent, $-(x - \mu)^2/(2\sigma^2)$, which is
minus one half of the squared number of standard deviations from the mean.

\begin{definition}[z-score]\index{z-score}
The \emph{z-score} of a value is its distance from the mean in units of standard deviation,
$z = (x - \mu)/\sigma$. If $x \sim \mathcal{N}(\mu, \sigma^2)$ then $z \sim \mathcal{N}(0, 1)$, the
\emph{standard normal}\index{standard normal}, with density
$\varphi(z) = \frac{1}{\sqrt{2\pi}}e^{-z^2/2}$.
\end{definition}

Substituting $z$ into \cref{eq:rx:gauss1d} shows the exponent is $-z^2/2$: the density of a value
depends on it only through $z^2$. Values with the same $z^2$ (the same distance from the mean, on
either side) are equally likely, and a larger $z^2$ means an exponentially less likely value. So
$z^2$ is a natural surprise score in one dimension. Remember that sentence: the multi-band score is
the same statement with $z^2$ replaced by its multi-band version.

\subsection{Tail probabilities and thresholds}

The probability that a standard normal lands more than $t$ from zero in either direction is the
area in its two tails, $P(\abs{z} > t)$. There is no closed form for it, but the values are
tabulated (and computed by \code{scipy.stats.norm.sf}):
\begin{center}
\begin{tabular}{ccc}
\toprule
$t$ & $P(\abs{z} > t)$ & roughly\\
\midrule
1    & 0.3173   & one pixel in 3\\
2    & 0.0455   & one in 22\\
3    & 0.0027   & one in 370\\
3.42 & 0.000626 & one in 1{,}600\\
\bottomrule
\end{tabular}
\end{center}
(\texttt{tools/compute\_rx.py} computes every number in this chapter.) If the thermal scene really
were Gaussian, a pixel at $z = 3.42$ or beyond would occur once in about 1,600 pixels; a Landsat
scene has millions of pixels, so you would still expect thousands of them. That is the first
lesson of detection: \emph{unusual} is relative to how many pixels you look at.

Turn the table around and you have a \emph{threshold}\index{threshold}. Decide what fraction of
normal pixels you are willing to flag by mistake, the \emph{false-alarm rate}\index{false-alarm
rate} $\alpha$, and choose $t$ so that $P(\abs{z} > t) = \alpha$. For $\alpha = 1\%$, $t = 2.576$;
for $\alpha = 0.1\%$, $t = 3.291$. In terms of the score $z^2$ the thresholds are
$2.576^2 = 6.635$ and $3.291^2 = 10.83$.

\begin{figure}[htbp]
\centering
\begin{tikzpicture}
\begin{groupplot}[group style={group size=2 by 1,horizontal sep=1.6cm},
  height=5.2cm,samples=121,axis lines=left,xlabel={$z$},
  tick label style={font=\scriptsize},label style={font=\small}]
\nextgroupplot[width=0.58\linewidth,domain=-4:4,ymin=0,ymax=0.43,
  ylabel={density $\varphi(z)$},xtick={-2.58,0,2.58},
  xticklabels={$-2.58$,0,$2.58$}]
  \addplot[fill=loss!35,draw=none,domain=2.576:4] {exp(-x^2/2)/sqrt(2*pi)} \closedcycle;
  \addplot[fill=loss!35,draw=none,domain=-4:-2.576] {exp(-x^2/2)/sqrt(2*pi)} \closedcycle;
  \addplot[fill=loss!85,draw=none,domain=3.291:4] {exp(-x^2/2)/sqrt(2*pi)} \closedcycle;
  \addplot[fill=loss!85,draw=none,domain=-4:-3.291] {exp(-x^2/2)/sqrt(2*pi)} \closedcycle;
  \addplot[spec,thick] {exp(-x^2/2)/sqrt(2*pi)};
  \node[font=\scriptsize,align=center] at (axis cs:0,0.17) {99\%\\ of normal\\ pixels};
  \draw[maskgrey] (axis cs:2.2,0) rectangle (axis cs:4,0.03);
\nextgroupplot[width=0.42\linewidth,domain=2.2:4,xmin=2.2,xmax=4,ymin=0,ymax=0.03,
  title={\scriptsize zoom on the right tail},xtick={2.58,3.29},
  xticklabels={2.58,3.29},yticklabel style={/pgf/number format/fixed}]
  \addplot[fill=loss!35,draw=none,domain=2.576:4] {exp(-x^2/2)/sqrt(2*pi)} \closedcycle;
  \addplot[fill=loss!85,draw=none,domain=3.291:4] {exp(-x^2/2)/sqrt(2*pi)} \closedcycle;
  \addplot[spec,thick] {exp(-x^2/2)/sqrt(2*pi)};
  \draw[attn,thick,dashed] (axis cs:3.42,0) -- (axis cs:3.42,0.022)
    node[above,font=\scriptsize,align=center] {thermal pixel\\$z = 3.42$};
\end{groupplot}
\end{tikzpicture}
\caption{The standard normal density. The light and dark tails together hold 1\% of the area
(beyond $\pm2.576$); the dark tails alone hold 0.1\% (beyond $\pm3.291$). They are thin, so the
right panel magnifies the grey box on the right tail. The thermal pixel of \cref{sec:rx:surprise}
sits at $z = 3.42$, beyond both thresholds.}
\label{fig:rx:normal}
\end{figure}

\section{Why many bands need more than many z-scores}\label{sec:rx:whynotz}

The obvious way to score a 165-band pixel is to compute a z-score per band and add the squares.
That fails because \textbf{bands are not independent}. Vegetation is dark in red and bright in the
near infrared, and across a scene the two bands move together: more vegetation, darker red and
brighter infrared. A pixel slightly bright in both is ordinary. A pixel bright in red \emph{and}
dark in the infrared is bizarre, even if neither value alone is extreme. Per-band z-scores look at
one band at a time and cannot see the relationship.

\begin{worked}[label={we:rx:zfail}]{Per-band z-scores miss the pattern}
Use the two-band covariance of \cref{we:linalg:cov}, $\Cov = \begin{psmallmatrix} 2 & 3\\ 3 &
10\end{psmallmatrix}$ (spreads $\sqrt2$ and $\sqrt{10}$, positive covariance). Two test pixels sit
these distances from the mean:
\begin{align*}
\vect{d}_A &= (2, 6) && \text{(both bands up, the way the scene moves)},\\
\vect{d}_B &= (2, -3) && \text{(band 1 up, band 2 down)}.
\end{align*}
Per-band squared z-scores, $d_1^2/2 + d_2^2/10$:
\begin{align*}
A:&\quad \frac{2^2}{2} + \frac{6^2}{10} = \frac42 + \frac{36}{10} = 2 + 3.6 = 5.6\\
B:&\quad \frac{2^2}{2} + \frac{(-3)^2}{10} = \frac42 + \frac{9}{10} = 2 + 0.9 = 2.9 .
\end{align*}
By this score A is nearly twice as unusual as B. Now use the covariance properly. With
$\Cov\inv = \frac{1}{11}\begin{psmallmatrix} 10 & -3\\ -3 & 2\end{psmallmatrix}$ (\cref{we:linalg:inv}):
\begin{align*}
\Cov\inv\vect{d}_A &= \tfrac{1}{11}(10\times2 - 3\times6,\ -3\times2 + 2\times6)
  = \tfrac{1}{11}(20 - 18,\ -6 + 12) = \tfrac{1}{11}(2, 6)\\
\vect{d}_A\T\Cov\inv\vect{d}_A &= \tfrac{1}{11}(2\times2 + 6\times6) = \tfrac{1}{11}(4 + 36)
  = \tfrac{40}{11} = 3.64\\[4pt]
\Cov\inv\vect{d}_B &= \tfrac{1}{11}(10\times2 - 3\times(-3),\ -3\times2 + 2\times(-3))
  = \tfrac{1}{11}(20 + 9,\ -6 - 6) = \tfrac{1}{11}(29, -12)\\
\vect{d}_B\T\Cov\inv\vect{d}_B &= \tfrac{1}{11}(2\times29 + (-3)\times(-12)) = \tfrac{1}{11}(58 + 36)
  = \tfrac{94}{11} = 8.55 .
\end{align*}
The ranking flips. A moves along the direction the scene always moves, and scores 3.64; B breaks the
relationship between the bands, and scores 8.55. The quantity $\vect{d}\T\Cov\inv\vect{d}$ is the
Mahalanobis distance, and the rest of the chapter explains why it is the right score.
\end{worked}

\section{The multivariate Gaussian}\label{sec:rx:mvn}

To score many bands the way $z^2$ scores one, you need a probability model for a whole spectrum: a
density over $\R^C$ that has a mean vector, a covariance matrix, and a bell shape. Build it from the
one-dimensional case in three steps.

\paragraph{Step 1: independent standard normals.} Let $\vect{z} = (z_1, \dots, z_C)$ have entries
that are independent standard normals. For independent quantities the joint density is the product
of the individual densities, and products of exponentials add their exponents:
\begin{align}
p(\vect{z}) &= \prod_{k=1}^{C}\frac{1}{\sqrt{2\pi}}e^{-z_k^2/2}
  = \frac{1}{(2\pi)^{C/2}}\exp\!\left(-\tfrac12\left(z_1^2 + \cdots + z_C^2\right)\right) \notag\\
&= \frac{1}{(2\pi)^{C/2}}\exp\!\left(-\tfrac12\norm{\vect{z}}^2\right). \label{eq:rx:stdmvn}
\end{align}
The density depends on $\vect{z}$ only through its length: the contours of equal density are
circles (spheres in more dimensions). Its mean is $\vect{0}$ and its covariance is $\mat{I}$ (each
entry has variance 1, and independent entries have covariance 0). It is exactly the round, whitened
cloud of \cref{fig:linalg:whitening}(b).

\paragraph{Step 2: stretch and turn it.} Take a mean $\muv$ and a covariance $\Cov =
\mat{U}\mat{\Lambda}\mat{U}\T$ with all eigenvalues positive. Define
\[
\vect{x} = \muv + \mat{A}\vect{z}, \qquad \mat{A} = \mat{U}\mat{\Lambda}^{1/2}.
\]
In words: stretch the round cloud by $\sqrt{\lambda_k}$ along axis $k$, rotate the axes onto the
eigenvectors, and move the centre to $\muv$. By \cref{eq:linalg:covtransform} the covariance of
$\vect{x}$ is $\mat{A}\mat{I}\mat{A}\T = \mat{U}\mat{\Lambda}^{1/2}\mat{\Lambda}^{1/2}\mat{U}\T =
\mat{U}\mat{\Lambda}\mat{U}\T = \Cov$, as wanted.

\paragraph{Step 3: change variables.} What is the density of $\vect{x}$? First the one-dimensional
version. If $x = \mu + \sigma z$ then $z = (x - \mu)/\sigma$, and an interval of width $\Delta x$ in
$x$ corresponds to an interval of width $\Delta x/\sigma$ in $z$. The same probability is spread
over an interval $\sigma$ times wider, so the density is $\sigma$ times lower:
$p_x(x) = \frac{1}{\sigma}\,\varphi\!\left(\frac{x - \mu}{\sigma}\right)$, which is exactly
\cref{eq:rx:gauss1d}. In $C$ dimensions the same argument uses volume instead of width: the map
$\vect{z} \mapsto \mat{A}\vect{z}$ multiplies volumes by $\abs{\det\mat{A}}$, so
\[
p_x(\vect{x}) = \frac{1}{\abs{\det\mat{A}}}\;p_z\!\left(\mat{A}\inv(\vect{x} - \muv)\right).
\]
Now identify the pieces. The inverse of $\mat{A} = \mat{U}\mat{\Lambda}^{1/2}$ is
$\mat{\Lambda}^{-1/2}\mat{U}\T$ (check: $\mat{\Lambda}^{-1/2}\mat{U}\T\mat{U}\mat{\Lambda}^{1/2} =
\mat{\Lambda}^{-1/2}\mat{\Lambda}^{1/2} = \mat{I}$), and that is the whitening matrix $\mat{W}$ of
\cref{eq:linalg:whiten}. The determinant of a product is the product of determinants, an orthogonal
matrix has determinant $\pm1$, and a diagonal matrix's determinant is the product of its diagonal,
so $\abs{\det\mat{A}} = \sqrt{\lambda_1\cdots\lambda_C} = \sqrt{\det\Cov}$. Substituting into
\cref{eq:rx:stdmvn}:

\begin{definition}[Multivariate Gaussian]\index{Gaussian distribution!multivariate}
A spectrum $\vect{x} \in \R^C$ has a multivariate Gaussian distribution with mean $\muv$ and
positive definite covariance $\Cov$, written $\vect{x}\sim\mathcal{N}(\muv, \Cov)$, if
\begin{equation}\label{eq:rx:mvn}
\begin{split}
p(\vect{x}) &= \frac{1}{(2\pi)^{C/2}\sqrt{\det\Cov}}
\exp\!\left(-\tfrac12\,\norm{\mat{W}(\vect{x} - \muv)}^2\right)\\
&= \frac{1}{(2\pi)^{C/2}\sqrt{\det\Cov}}
\exp\!\left(-\tfrac12\,(\vect{x} - \muv)\T\Cov\inv(\vect{x} - \muv)\right).
\end{split}
\end{equation}
\end{definition}

The second form uses the identity $\norm{\mat{W}\vect{d}}^2 = \vect{d}\T\Cov\inv\vect{d}$, which the
next section proves. With $C = 1$, $\Cov = \sigma^2$, and \cref{eq:rx:mvn} is \cref{eq:rx:gauss1d}.
As in one dimension, everything that depends on the pixel sits in the exponent, and the exponent
is minus one half of a squared distance.

\section{Mahalanobis distance}\label{sec:rx:mahalanobis}

\begin{definition}[Mahalanobis distance]\index{Mahalanobis distance}
The (squared) \emph{Mahalanobis distance} of a vector $\vect{x}$ from a distribution with mean
$\muv$ and covariance $\Cov$ is
\begin{equation}\label{eq:rx:maha}
D^2(\vect{x}) = (\vect{x} - \muv)\T\Cov\inv(\vect{x} - \muv).
\end{equation}
\end{definition}

It is distance from the mean measured in units of how much the data actually vary in that
direction. Ten units along a direction in which the scene varies a lot are unremarkable; ten units
along a direction in which it hardly varies at all are extraordinary. The next result makes that
sentence exact.

\subsection{Mahalanobis distance is Euclidean distance after whitening}

\paragraph{Claim.} With $\Cov = \mat{U}\mat{\Lambda}\mat{U}\T$ and
$\mat{W} = \mat{\Lambda}^{-1/2}\mat{U}\T$,
\begin{equation}\label{eq:rx:mahawhite}
(\vect{x} - \muv)\T\Cov\inv(\vect{x} - \muv) = \norm{\mat{W}(\vect{x} - \muv)}^2
= \sum_{k=1}^{C}\frac{\left(\vect{u}_k\T(\vect{x} - \muv)\right)^2}{\lambda_k}.
\end{equation}

\noindent\textit{Proof.} Three steps.

\emph{(i) The inverse of the covariance.} Guess $\Cov\inv = \mat{U}\mat{\Lambda}\inv\mat{U}\T$, with
$\mat{\Lambda}\inv = \diag(1/\lambda_1, \dots, 1/\lambda_C)$, and check by multiplying:
\begin{align*}
\Cov\,\mat{U}\mat{\Lambda}\inv\mat{U}\T
&= \mat{U}\mat{\Lambda}\mat{U}\T\mat{U}\mat{\Lambda}\inv\mat{U}\T
  && \text{spectral theorem}\\
&= \mat{U}\mat{\Lambda}\mat{\Lambda}\inv\mat{U}\T
  && \mat{U}\T\mat{U} = \mat{I}\\
&= \mat{U}\mat{U}\T = \mat{I}
  && \text{diagonal entries } \lambda_k/\lambda_k = 1;\ \mat{U}\text{ is square and orthogonal.}
\end{align*}
So the inverse has the same eigenvectors and the reciprocal eigenvalues. Remember this: \emph{a
small eigenvalue of $\Cov$ is a large eigenvalue of $\Cov\inv$.}

\emph{(ii) $\mat{W}\T\mat{W} = \Cov\inv$.} Using \cref{eq:linalg:abt},
\[
\mat{W}\T\mat{W} = (\mat{\Lambda}^{-1/2}\mat{U}\T)\T\mat{\Lambda}^{-1/2}\mat{U}\T
= \mat{U}\mat{\Lambda}^{-1/2}\mat{\Lambda}^{-1/2}\mat{U}\T
= \mat{U}\mat{\Lambda}\inv\mat{U}\T = \Cov\inv .
\]

\emph{(iii) Put them together.} Write $\vect{d} = \vect{x} - \muv$. Then
\[
\norm{\mat{W}\vect{d}}^2 = (\mat{W}\vect{d})\T(\mat{W}\vect{d}) = \vect{d}\T\mat{W}\T\mat{W}\vect{d}
= \vect{d}\T\Cov\inv\vect{d}.
\]
For the sum form, entry $k$ of $\mat{W}\vect{d}$ is row $k$ of $\mat{W}$ times $\vect{d}$, which is
$\vect{u}_k\T\vect{d}/\sqrt{\lambda_k}$; squaring and adding gives the right-hand side of
\cref{eq:rx:mahawhite}. \hfill$\square$

\begin{keypoint}
Mahalanobis distance = turn the pixel onto the principal directions, divide each coordinate by its
standard deviation, and take the ordinary squared length.
\end{keypoint}

The sum form is the one to hold in your head. Each principal direction contributes its squared
coordinate divided by its eigenvalue. A direction with a large eigenvalue (the scene varies a lot
that way) contributes little; a direction with a small eigenvalue contributes a lot. The contours of
equal Mahalanobis distance are the ellipses of \cref{fig:linalg:whitening}(a), long along the large
eigenvectors and short along the small ones.

\begin{worked}[label={we:rx:twoways}]{Two ways to the same Mahalanobis distance}
Take pixel B of \cref{we:rx:zfail}, $\vect{d}_B = (2, -3)$, with $\vect{u}_1 = (1, 3)/\sqrt{10}$,
$\lambda_1 = 11$, $\vect{u}_2 = (3, -1)/\sqrt{10}$, $\lambda_2 = 1$ (\cref{we:linalg:eig}).
Coordinates:
\begin{align*}
\vect{u}_1\T\vect{d}_B &= \frac{1\times2 + 3\times(-3)}{\sqrt{10}} = \frac{2 - 9}{\sqrt{10}}
  = \frac{-7}{\sqrt{10}}\\
\vect{u}_2\T\vect{d}_B &= \frac{3\times2 + (-1)\times(-3)}{\sqrt{10}} = \frac{6 + 3}{\sqrt{10}}
  = \frac{9}{\sqrt{10}} .
\end{align*}
Squared, divided by the eigenvalues, added:
\[
D^2 = \frac{49/10}{11} + \frac{81/10}{1} = \frac{49}{110} + \frac{81}{10}
= 0.445 + 8.100 = 8.545 = \frac{94}{11}.
\]
The same 8.55 as \cref{we:rx:zfail}, and now you can see where it comes from: 8.1 of it is the
pixel's step of $9/\sqrt{10} = 2.85$ along $\vect{u}_2$, the direction in which the scene's spread is
only 1.
\end{worked}

\section{The RX detector}\label{sec:rx:rx}

The \emph{RX detector}\index{RX detector}, named after Reed and Xiaoli Yu, who derived it in 1990
\citep{reed1990adaptive}, scores every pixel by its Mahalanobis distance from the scene's background
statistics:
\begin{equation}\label{eq:rx:rx}
\text{RX}(\vect{x}) = (\vect{x} - \hat{\muv})\T\hat{\Cov}\inv(\vect{x} - \hat{\muv}),
\end{equation}
where the hats mean ``estimated from the pixels'': $\hat\muv$ is the scene's mean spectrum and
$\hat\Cov$ its covariance matrix (\cref{eq:linalg:mean,eq:linalg:cov}). \emph{Global RX} estimates
them once from all valid pixels of the scene; \emph{local RX} (\cref{ch:localrx}) estimates them
separately for each pixel from a ring of its neighbours.

\subsection{Where the formula comes from: a likelihood ratio}
\index{generalised likelihood ratio}

Why this score and not some other? Because it is what a likelihood-ratio test gives. Frame the
decision at one pixel as two hypotheses:
\begin{itemize}
  \item $H_0$ (background): $\vect{x} \sim \mathcal{N}(\muv, \Cov)$;
  \item $H_1$ (target): $\vect{x} \sim \mathcal{N}(\muv + \vect{s}, \Cov)$ for some unknown signature
    shift $\vect{s}$.
\end{itemize}
A \emph{likelihood ratio}\index{likelihood ratio} compares how well each hypothesis explains the
observed pixel: $p(\vect{x}\mid H_1)/p(\vect{x}\mid H_0)$. We do not know $\vect{s}$, so the
\emph{generalised} likelihood ratio uses the $\vect{s}$ that explains the pixel best. Under $H_1$
the exponent is $-\frac12(\vect{x} - \muv - \vect{s})\T\Cov\inv(\vect{x} - \muv - \vect{s})$, which is
at most 0 and equals 0 when $\vect{s} = \vect{x} - \muv$. So the best $H_1$ density is the peak
value $1/((2\pi)^{C/2}\sqrt{\det\Cov})$, and taking the logarithm of the ratio:
\begin{align*}
\log\frac{\max_{\vect{s}} p(\vect{x}\mid H_1)}{p(\vect{x}\mid H_0)}
&= \log\frac{\frac{1}{(2\pi)^{C/2}\sqrt{\det\Cov}}\,e^{0}}
  {\frac{1}{(2\pi)^{C/2}\sqrt{\det\Cov}}\,e^{-\frac12 D^2(\vect{x})}}
  && \text{\cref{eq:rx:mvn} for both}\\
&= \log e^{\frac12 D^2(\vect{x})}
  && \text{the normalising factors cancel}\\
&= \tfrac12\,D^2(\vect{x}) .
\end{align*}
The log likelihood ratio is half the Mahalanobis distance, and a larger ratio means stronger
evidence for a target. Ranking by the ratio and ranking by $D^2$ are the same ranking.

\begin{sourcenote}[: the original RX statistic]
Reed and Yu derived their test with the background mean and covariance themselves estimated from a
finite set of $N$ background pixels \citep{reed1990adaptive}, which adds terms involving $N$ to the
statistic; it is still an increasing function of the Mahalanobis distance computed with the
estimates, so it ranks pixels the same way, and with the large backgrounds used here it reduces to
\cref{eq:rx:rx}. The code computes \cref{eq:rx:rx} directly: \code{GlobalRXDetector.detect}
(\code{app/detectors/global_rx_detector.py}, lines 227--229) calls \code{spectral.rx} on the valid
pixels, which estimates $\hat\muv$ and $\hat\Cov$ from those same pixels (the \code{spectral}
library's \code{calc_stats}, divisor $N - 1$) and forms $\hat\Cov\inv$ with \code{np.linalg.inv}.
\end{sourcenote}

Notice what the global version implies: the statistics are estimated from \emph{all} valid pixels,
including any anomalies. A handful of anomalous pixels among millions barely moves the mean and
covariance, so for global RX this contamination is mild. For local RX, where the background is a
few thousand pixels around the target, it is not mild at all, and \cref{ch:localrx} builds a guard
window against it.

\section{How large is large? The chi-square distribution}\label{sec:rx:chi2}

If the background really is Gaussian, the RX score of a background pixel has a known distribution,
and that tells you what a ``large'' score is.

By \cref{eq:rx:mahawhite}, $D^2 = \norm{\vect{z}}^2 = z_1^2 + \cdots + z_C^2$ with
$\vect{z} = \mat{W}(\vect{x} - \muv)$. If $\vect{x} \sim \mathcal{N}(\muv, \Cov)$ then, running the
construction of \cref{sec:rx:mvn} backwards, $\vect{z}$ has independent standard normal entries. So
under the Gaussian background,
\begin{equation}\label{eq:rx:chi2sum}
D^2 = z_1^2 + z_2^2 + \cdots + z_C^2,\qquad z_k \text{ independent } \mathcal{N}(0, 1).
\end{equation}

\begin{definition}[Chi-square distribution]\index{chi-square distribution}
The distribution of a sum of $C$ squared independent standard normals is the \emph{chi-square
distribution with $C$ degrees of freedom}, written $\chi^2_C$. Its density, for $t > 0$, is
\begin{equation}\label{eq:rx:chi2pdf}
f_C(t) = \frac{t^{C/2 - 1}\,e^{-t/2}}{2^{C/2}\,\Gamma(C/2)},
\end{equation}
where $\Gamma$ is the gamma function ($\Gamma(n) = (n-1)!$ for whole numbers $n$,
$\Gamma(1/2) = \sqrt\pi$, $\Gamma(3/2) = \sqrt\pi/2$).
\end{definition}

You will not need to derive \cref{eq:rx:chi2pdf} in general, but three facts about it you should be
able to derive yourself.

\paragraph{The mean is $C$.} Each $z_k^2$ has mean $\E[z_k^2] = \operatorname{Var}(z_k) +
(\E z_k)^2 = 1 + 0 = 1$, and means add, so $\E[D^2] = C$. A background pixel in a 10-dimensional
space has an average RX score of 10, not 1. Scores grow with the number of bands even when nothing
is anomalous, which is why you can never compare RX scores computed in different dimensions
directly.

\paragraph{The spread is $\sqrt{2C}$.} For a standard normal $\E[z^4] = 3$ (a standard fact), so
$\operatorname{Var}(z^2) = \E[z^4] - (\E[z^2])^2 = 3 - 1 = 2$. Variances of independent quantities
add, so $\operatorname{Var}(D^2) = 2C$ and the standard deviation is $\sqrt{2C}$: 4.47 for
$C = 10$, 18.2 for $C = 165$.

\paragraph{Two cases by hand.} For $C = 1$, $D^2 = z^2$ and $P(z^2 > t) = P(\abs{z} > \sqrt t)$: the
one-dimensional tail table of \cref{sec:rx:gauss1d}, with $t$ squared. For $C = 2$ the tail has a
closed form. The density of $(z_1, z_2)$ is $\frac{1}{2\pi}e^{-r^2/2}$ with $r^2 = z_1^2 + z_2^2$;
integrate over the region outside the circle of radius $\sqrt t$ in polar coordinates (the ring at
radius $r$ has circumference $2\pi r$):
\[
P(D^2 > t) = \int_{\sqrt t}^{\infty}\frac{1}{2\pi}e^{-r^2/2}\,2\pi r\,dr
= \Big[-e^{-r^2/2}\Big]_{\sqrt t}^{\infty} = 0 - (-e^{-t/2}) = e^{-t/2}.
\]
So for two bands the threshold with false-alarm rate $\alpha$ solves $e^{-t/2} = \alpha$:
$t = 2\ln(1/\alpha)$. For $\alpha = 1\%$, $t = 2\ln 100 = 2\times4.605 = 9.21$.

\begin{figure}[htbp]
\centering
\begin{tikzpicture}
\begin{groupplot}[group style={group size=2 by 1,horizontal sep=1.5cm},
  height=5.6cm,axis lines=left,xlabel={score $t = D^2$},
  tick label style={font=\scriptsize},label style={font=\small}]
\nextgroupplot[width=0.56\linewidth,xmin=0,xmax=32,ymin=0,ymax=0.5,ylabel={density $f_C(t)$},
  legend style={font=\scriptsize,draw=none,at={(0.98,0.97)},anchor=north east}]
  \addplot[spec,thick,domain=0.05:32,samples=200] {exp(-x/2)/sqrt(2*pi*x)};
  \addlegendentry{$C = 1$}
  \addplot[attn,thick,domain=0:32,samples=120] {0.5*exp(-x/2)};
  \addlegendentry{$C = 2$}
  \addplot[spat,thick,domain=0.001:32,samples=200] {sqrt(x)*exp(-x/2)/sqrt(2*pi)};
  \addlegendentry{$C = 3$}
  \addplot[loss,thick,domain=0:32,samples=200] {x^4*exp(-x/2)/768};
  \addlegendentry{$C = 10$}
  \draw[loss,dotted] (axis cs:10,0) -- (axis cs:10,0.14) node[above,font=\scriptsize] {mean 10};
  \draw[maskgrey] (axis cs:8,0) rectangle (axis cs:32,0.04);
\nextgroupplot[width=0.44\linewidth,xmin=8,xmax=32,ymin=0,ymax=0.04,
  title={\scriptsize zoom: the 1\% tails},yticklabel style={/pgf/number format/fixed},
  xtick={11.34,23.21},xticklabels={{\color{spat}11.34},{\color{loss}23.21}}]
  \addplot[fill=spat!35,draw=none,domain=11.34:32,samples=100] {sqrt(x)*exp(-x/2)/sqrt(2*pi)} \closedcycle;
  \addplot[fill=loss!40,draw=none,domain=23.21:32,samples=100] {x^4*exp(-x/2)/768} \closedcycle;
  \addplot[spat,thick,domain=8:32,samples=150] {sqrt(x)*exp(-x/2)/sqrt(2*pi)};
  \addplot[loss,thick,domain=8:32,samples=150] {x^4*exp(-x/2)/768};
\end{groupplot}
\end{tikzpicture}
\caption{Chi-square densities for $C = 1, 2, 3, 10$ degrees of freedom. The right panel magnifies the
grey box: its shaded tails hold 1\% of the area for $C = 3$ (beyond 11.34) and for $C = 10$ (beyond
23.21). As $C$ grows the whole
distribution moves right (mean $C$) and spreads (standard deviation $\sqrt{2C}$), so a score of 15
is alarming in 3 dimensions and ordinary in 10.}
\label{fig:rx:chi2}
\end{figure}

\subsection{Thresholds by false-alarm rate}
\index{chi-square distribution!thresholds}

The $\chi^2_C$ tail gives a threshold for any false-alarm rate. The values below come from
\code{scipy.stats.chi2.isf}\tagC:
\begin{center}
\begin{tabular}{rcccc}
\toprule
$C$ & mean & standard deviation & threshold at $\alpha = 1\%$ & at $\alpha = 0.1\%$\\
\midrule
1   & 1   & 1.41  & 6.63   & 10.83\\
2   & 2   & 2.00  & 9.21   & 13.82\\
3   & 3   & 2.45  & 11.34  & 16.27\\
10  & 10  & 4.47  & 23.21  & 29.59\\
30  & 30  & 7.75  & 50.89  & 59.70\\
165 & 165 & 18.17 & 210.18 & 226.88\\
\bottomrule
\end{tabular}
\end{center}
The $C = 1$ row matches the squared normal thresholds of \cref{sec:rx:gauss1d} (6.635 and 10.83), and
the $C = 2$ row matches $2\ln100$ and $2\ln1000$. The $C = 10$ row is the one that would apply to
MNF-RX with its 10 components (\cref{ch:localrx}), and $C = 30$ to the 30-component scores of
\cref{ch:localrx,ch:residualmaha}.

\begin{keypoint}
Under a Gaussian background the RX score of a normal pixel is $\chi^2_C$: average $C$, and a
threshold for any false-alarm rate can be read from a table.
\end{keypoint}

\paragraph{Why the product does not use these thresholds.} Real scenes are not one Gaussian. They
are mixtures of surfaces, with heavy tails from edges, glint and stripes, so the fraction of real
pixels above the $\chi^2$ threshold is usually far more than $\alpha$. Every score in the product is
therefore thresholded by \emph{percentile} of its own map (the human threshold step of
\cref{ch:product}), and drawn with a square-root stretch because squared distances have a long right
tail. The chi-square result is still worth knowing: it tells you what scale to expect (scores near
$C$ are ordinary), it explains why scores computed in different numbers of dimensions are not
comparable, and it is the reason the residual score's docstring calls its output
``chi-squared-like''.

\subsection{A fact that holds for any data}

One property of RX scores does not need the Gaussian at all. When the mean and covariance are
estimated from the same $N$ pixels that are scored (as \code{spectral.rx} does), the scores always
average to exactly $C(N-1)/N$. The proof uses the \emph{trace}\index{trace}, the sum of a square
matrix's diagonal, and two of its properties: it is linear, and for vectors $\vect{a}, \vect{b}$,
$\vect{a}\T\vect{b} = \tr(\vect{b}\vect{a}\T)$ (both are $\sum_i a_ib_i$; the second is the sum of the
diagonal of the outer product). Then
\begin{align*}
\sum_{n=1}^{N}\vect{d}_n\T\hat\Cov\inv\vect{d}_n
&= \sum_n\tr\!\left(\hat\Cov\inv\vect{d}_n\vect{d}_n\T\right)
  && \vect{a} = \vect{d}_n,\ \vect{b} = \hat\Cov\inv\vect{d}_n\\
&= \tr\!\Big(\hat\Cov\inv\sum_n\vect{d}_n\vect{d}_n\T\Big)
  && \text{the trace and the matrix product are linear}\\
&= \tr\!\left(\hat\Cov\inv(N - 1)\hat\Cov\right)
  && \text{\cref{eq:linalg:cov}}\\
&= (N - 1)\tr(\mat{I}_C) = (N - 1)\,C .
\end{align*}
Dividing by $N$ gives the average. This is a useful sanity check on any implementation: the mean
in-sample score of a correct RX is $C(N-1)/N$, whatever the data look like. You will verify it on
the worked example below.

\section{A full worked example in three bands}\label{sec:rx:worked3}

\begin{sourcenote}[: the course's three-band table cannot be scored]
The anomaly course builds a covariance from four pixels in three bands, $(0.10, 0.20, 0.30)$,
$(0.12, 0.24, 0.28)$, $(0.08, 0.16, 0.32)$ and $(0.10, 0.20, 0.30)$
(\file{docs/training/anomaly-detection/08-rx-detection.md}, ``Covariance, in three bands''). Its
deviations are $\vect{0}$, $(0.02, 0.04, -0.02)$, $-(0.02, 0.04, -0.02)$ and $\vect{0}$: every one is a
multiple of the same vector. Their covariance is therefore $\frac{0.0008}{3}\,\vect{v}\vect{v}\T$
with $\vect{v} = (1, 2, -1)$, a matrix of rank 1 with eigenvalues $0.0016$, $0$ and $0$
(\texttt{tools/compute\_rx\_output.txt}). Its determinant is 0 and it has no inverse, so RX cannot
score anything against it. The table is a fine illustration of covariance signs, which is what the
course uses it for, and an accidental illustration of the singular covariance of
\cref{sec:rx:singular}. The example below uses five pixels that span all three directions.
\end{sourcenote}

\begin{worked}[label={we:rx:three}]{RX by hand on three bands}
Five background pixels, three bands, reflectance in percent. This is the longest calculation in the
chapter; do it on paper, one step at a time.
\begin{center}
\begin{tabular}{cccc}
\toprule
pixel & band 1 & band 2 & band 3\\
\midrule
1 & 8  & 18 & 30\\
2 & 9  & 21 & 30\\
3 & 11 & 19 & 28\\
4 & 11 & 21 & 30\\
5 & 11 & 21 & 32\\
\bottomrule
\end{tabular}
\end{center}
\textbf{Step 1, the mean.}
\begin{align*}
\mu_1 &= (8 + 9 + 11 + 11 + 11)/5 = 50/5 = 10\\
\mu_2 &= (18 + 21 + 19 + 21 + 21)/5 = 100/5 = 20\\
\mu_3 &= (30 + 30 + 28 + 30 + 32)/5 = 150/5 = 30 .
\end{align*}
\textbf{Step 2, the deviations} $\vect{d}_n = \vect{x}_n - (10, 20, 30)$:
\[
\vect{d}_1 = (-2, -2, 0),\quad \vect{d}_2 = (-1, 1, 0),\quad \vect{d}_3 = (1, -1, -2),\quad
\vect{d}_4 = (1, 1, 0),\quad \vect{d}_5 = (1, 1, 2).
\]
Each band's deviations add to zero: $-2 - 1 + 1 + 1 + 1 = 0$; $-2 + 1 - 1 + 1 + 1 = 0$;
$0 + 0 - 2 + 0 + 2 = 0$.

\textbf{Step 3, the covariance.} The six distinct entries are sums of products of deviations,
divided by $N - 1 = 4$:
\begin{align*}
\Sigma_{11} &= \tfrac14\left[(-2)^2 + (-1)^2 + 1^2 + 1^2 + 1^2\right] = \tfrac14(4 + 1 + 1 + 1 + 1)
  = \tfrac84 = 2\\
\Sigma_{22} &= \tfrac14\left[(-2)^2 + 1^2 + (-1)^2 + 1^2 + 1^2\right] = \tfrac84 = 2\\
\Sigma_{33} &= \tfrac14\left[0 + 0 + (-2)^2 + 0 + 2^2\right] = \tfrac84 = 2\\
\Sigma_{12} &= \tfrac14\left[(-2)(-2) + (-1)(1) + (1)(-1) + (1)(1) + (1)(1)\right]
  = \tfrac14(4 - 1 - 1 + 1 + 1) = \tfrac44 = 1\\
\Sigma_{23} &= \tfrac14\left[(-2)(0) + (1)(0) + (-1)(-2) + (1)(0) + (1)(2)\right]
  = \tfrac14(0 + 0 + 2 + 0 + 2) = 1\\
\Sigma_{13} &= \tfrac14\left[(-2)(0) + (-1)(0) + (1)(-2) + (1)(0) + (1)(2)\right]
  = \tfrac14(0 + 0 - 2 + 0 + 2) = 0 .
\end{align*}
So
\[
\Cov = \begin{pmatrix} 2 & 1 & 0\\ 1 & 2 & 1\\ 0 & 1 & 2\end{pmatrix}.
\]
Each band has variance 2; neighbouring bands have covariance 1 (correlation $1/2$); bands 1 and 3
are uncorrelated. This is the pattern of a real spectrum in miniature: neighbours move together,
distant bands less so.

\textbf{Step 4, the determinant} by expansion along the first row (each entry times the determinant
of the $2\times2$ block left when you delete its row and column, with signs $+,-,+$):
\begin{align*}
\det\Cov &= 2\det\begin{pmatrix} 2 & 1\\ 1 & 2\end{pmatrix}
  - 1\det\begin{pmatrix} 1 & 1\\ 0 & 2\end{pmatrix}
  + 0\det\begin{pmatrix} 1 & 2\\ 0 & 1\end{pmatrix}\\
&= 2(4 - 1) - 1(2 - 0) + 0 = 6 - 2 = 4 .
\end{align*}
\textbf{Step 5, the inverse.}\index{cofactor}\index{inverse!of a 3x3 matrix} For a $3\times3$ matrix the recipe that generalises
\cref{eq:linalg:inv2} is: for each entry, compute its \emph{cofactor} (the determinant of the block
left after deleting its row and column, with sign $(-1)^{i+j}$), arrange the cofactors as a matrix,
transpose it, and divide by the determinant. Because $\Cov$ is symmetric, the cofactor matrix is
symmetric and the transpose changes nothing. The six distinct cofactors:
\begin{align*}
C_{11} &= +\det\begin{psmallmatrix} 2 & 1\\ 1 & 2\end{psmallmatrix} = 4 - 1 = 3, &
C_{12} &= -\det\begin{psmallmatrix} 1 & 1\\ 0 & 2\end{psmallmatrix} = -(2 - 0) = -2,\\
C_{13} &= +\det\begin{psmallmatrix} 1 & 2\\ 0 & 1\end{psmallmatrix} = 1 - 0 = 1, &
C_{22} &= +\det\begin{psmallmatrix} 2 & 0\\ 0 & 2\end{psmallmatrix} = 4 - 0 = 4,\\
C_{23} &= -\det\begin{psmallmatrix} 2 & 1\\ 0 & 1\end{psmallmatrix} = -(2 - 0) = -2, &
C_{33} &= +\det\begin{psmallmatrix} 2 & 1\\ 1 & 2\end{psmallmatrix} = 4 - 1 = 3 .
\end{align*}
So
\[
\Cov\inv = \frac14\begin{pmatrix} 3 & -2 & 1\\ -2 & 4 & -2\\ 1 & -2 & 3\end{pmatrix}.
\]
Check one row: row 1 of $\Cov$ times column 1 of the bracket is $2\times3 + 1\times(-2) + 0\times1 =
6 - 2 = 4$, and $4/4 = 1$; row 1 times column 2 is $2\times(-2) + 1\times4 + 0\times(-2) = 0$. Check
the rest yourself.

\textbf{Step 6, score a test pixel} $\vect{x} = (12, 19, 32)$. Its deviation is
$\vect{d} = (12 - 10, 19 - 20, 32 - 30) = (2, -1, 2)$. Multiply by the inverse, row by row:
\begin{align*}
\text{row 1:}&\quad \tfrac14\left[3\times2 + (-2)\times(-1) + 1\times2\right] = \tfrac14(6 + 2 + 2)
  = \tfrac{10}{4} = 2.5\\
\text{row 2:}&\quad \tfrac14\left[(-2)\times2 + 4\times(-1) + (-2)\times2\right] = \tfrac14(-4 - 4 - 4)
  = \tfrac{-12}{4} = -3\\
\text{row 3:}&\quad \tfrac14\left[1\times2 + (-2)\times(-1) + 3\times2\right] = \tfrac14(2 + 2 + 6)
  = \tfrac{10}{4} = 2.5 .
\end{align*}
Then the dot product with $\vect{d}$:
\[
D^2 = 2\times2.5 + (-1)\times(-3) + 2\times2.5 = 5 + 3 + 5 = 13 .
\]
\textbf{Step 7, compare.} The per-band score would be $\frac{2^2}{2} + \frac{(-1)^2}{2} +
\frac{2^2}{2} = 2 + 0.5 + 2 = 4.5$, about the mean of $\chi^2_3$ (3): unremarkable. The
Mahalanobis score is 13, beyond the 1\% threshold of 11.34 and with a $\chi^2_3$ tail probability of
0.0046. The pixel is high in bands 1 and 3 and low in band 2, between them, while the background's
neighbouring bands move together.

Now a second test pixel, $\vect{x}' = (12, 23, 32)$, with $\vect{d}' = (2, 3, 2)$: further from the
mean in every band, but in step with the background's correlations.
\begin{align*}
\Cov\inv\vect{d}' &= \tfrac14\left(6 - 6 + 2,\ -4 + 12 - 4,\ 2 - 6 + 6\right) = \tfrac14(2, 4, 2)
  = (0.5, 1, 0.5)\\
D^2 &= 2\times0.5 + 3\times1 + 2\times0.5 = 1 + 3 + 1 = 5 .
\end{align*}
Per-band it scores $2 + 4.5 + 2 = 8.5$, nearly twice the first pixel; by Mahalanobis it scores 5, an
ordinary value for $\chi^2_3$ (tail probability 0.17). The \code{spectral} library, given the five
background pixels' statistics, returns exactly $[13, 5]$\tagC.

\textbf{Step 8, check the in-sample identity.} Scoring the five background pixels against their own
statistics (\code{spectral.rx} on them) gives 3, 2.75, 2.75, 0.75 and 2.75\tagC. They add to
$3 + 2.75 + 2.75 + 0.75 + 2.75 = 12 = (N - 1)C = 4\times3$, as \cref{sec:rx:chi2} promised.
\end{worked}

\begin{worked}[label={we:rx:threeeig}]{The same score through the eigenvectors}
The eigenvalues of this $\Cov$ are $2 + \sqrt2 = 3.414$, $2$ and $2 - \sqrt2 = 0.586$, with unit
eigenvectors $\vect{u}_1 = (\tfrac12, \tfrac{1}{\sqrt2}, \tfrac12)$, $\vect{u}_2 = (\tfrac{1}{\sqrt2},
0, -\tfrac{1}{\sqrt2})$, $\vect{u}_3 = (\tfrac12, -\tfrac{1}{\sqrt2}, \tfrac12)$. (Check one:
$\Cov\vect{u}_3 = (1 - 0.7071,\ 0.5 - 1.4142 + 0.5,\ -0.7071 + 1) = (0.2929, -0.4142, 0.2929)$, and
$0.586\times(0.5, -0.7071, 0.5) = (0.2929, -0.4142, 0.2929)$.) The first direction is ``all bands
up together''; the third is ``middle band against the outer two''.

For $\vect{d} = (2, -1, 2)$ the coordinates are
\begin{align*}
\vect{u}_1\T\vect{d} &= 1 - 0.7071 + 1 = 1.2929,\\
\vect{u}_2\T\vect{d} &= 1.4142 - 1.4142 = 0,\\
\vect{u}_3\T\vect{d} &= 1 + 0.7071 + 1 = 2.7071,
\end{align*}
and \cref{eq:rx:mahawhite} gives
\[
D^2 = \frac{1.2929^2}{3.414} + \frac{0^2}{2} + \frac{2.7071^2}{0.586}
= \frac{1.6716}{3.414} + 0 + \frac{7.3284}{0.586} = 0.490 + 0 + 12.51 = 13.0 .
\]
Almost all of the score, 12.5 of 13, comes from the direction with the smallest eigenvalue. That is
the mechanism of RX in one line, and in the next section it becomes the problem.
\end{worked}

\section{Why global RX fails on real scenes}\label{sec:rx:fails}

Global RX makes two assumptions that a real hyperspectral scene breaks: that the background is one
Gaussian cloud, and that its covariance can be inverted safely.

\subsection{Mixed backgrounds}
\index{mixed background}\index{RX detector!global}

A scene is half cool forest and half warm bare ground. A hot patch of machinery sits on the bare
ground. Its temperature is high, but the scene-wide mean sits between forest and bare ground, and the
standard deviation is inflated by the split.

\begin{worked}[label={we:rx:mixture}]{A mixture hides an anomaly from the global mean}
Assume for illustration\tagA{} that the forest is at 20\,\textdegree C with standard deviation 2, the
bare ground at 40\,\textdegree C with standard deviation 2, half the scene each. The scene mean is
$0.5\times20 + 0.5\times40 = 30$. The scene variance is the average squared value minus the squared
mean; for each half the average squared value is its variance plus its squared mean:
\begin{align*}
\text{average of } x^2 &= 0.5\,(2^2 + 20^2) + 0.5\,(2^2 + 40^2) = 0.5\times404 + 0.5\times1604
  = 202 + 802 = 1004\\
\text{variance} &= 1004 - 30^2 = 1004 - 900 = 104,\qquad \text{standard deviation} = \sqrt{104} = 10.2 .
\end{align*}
The machinery is at 46\,\textdegree C. Against the whole scene, $z = (46 - 30)/10.2 = 1.57$: about
one pixel in nine of a Gaussian scene is that far out. Against the bare ground around it,
$z = (46 - 40)/2 = 3.0$. The global statistics hide it; the local ones reveal it.
\end{worked}

The global mean and covariance describe a cloud that no pixel belongs to: the average of two
separate clouds, stretched along the line between them. \Cref{ch:localrx} answers this with a local
background per pixel. Notice that the reconstruction family answers it differently: a network that
has learnt what each surface looks like leaves residuals that are much closer to one cloud, which is
why the residual score of \cref{ch:residualmaha} can use global statistics with less harm.

\subsection{The near-singular covariance}\label{sec:rx:singular}

This is the reason plain RX was taken out of the product.

\begin{sourcenote}[: why plain RX is not shipped]
\code{backend/allotrope/foundation_models/resolver.py}, lines 185--189:
``Plain RX on hyperspectral is intentionally not shipped (dropped 2026-05-11). With 150+ surviving
bands the covariance comes out near-singular and Mahalanobis distances blow up to 1e11. Use MNF-RX
instead -- same machinery, but RX runs in a 10-dim noise-whitened subspace so the covariance is
well-conditioned.'' The product's classical options are \code{mnf_rx} and \code{thermal_grx}. The
class \code{GlobalRXDetector} still exists and is registered under \code{ADModel.RX} in
\code{app/utils/anomaly_detection/detector_factory.py}, for notebooks and thermal use, but no
product model resolves to it. Do not reintroduce it on hyperspectral data.
\end{sourcenote}

A matrix is \emph{singular}\index{singular matrix} when its determinant is zero, which for a
covariance means some eigenvalue is zero: the pixels do not vary at all along that direction. It is
\emph{near-singular} when some eigenvalue is tiny compared with the largest. The size of the problem
is measured by the \emph{condition number}\index{condition number}
\begin{equation}\label{eq:rx:kappa}
\kappa(\Cov) = \frac{\lambda_{\max}}{\lambda_{\min}} .
\end{equation}
By step (i) of the proof in \cref{sec:rx:mahalanobis}, the eigenvalues of $\Cov\inv$ are
$1/\lambda_k$, so the largest entry-scale of $\Cov\inv$ is $1/\lambda_{\min}$: the inverse amplifies
the least-varying direction the most.

\begin{worked}[label={we:rx:nearsingular}]{Two highly correlated bands}
Two neighbouring bands, each with variance 1, correlation $\rho = 0.999$:
$\Cov = \begin{psmallmatrix} 1 & 0.999\\ 0.999 & 1\end{psmallmatrix}$.

\textbf{Eigenvalues.} Trace 2, determinant $1 - 0.999^2 = 1 - 0.998001 = 0.001999$. With $a = c = 1$
and $b = \rho$, \cref{eq:linalg:eig2} gives $\lambda = 1 \pm \rho$:
\[
\lambda_1 = 1.999 \text{ along } \vect{u}_1 = (1, 1)/\sqrt2,\qquad
\lambda_2 = 0.001 \text{ along } \vect{u}_2 = (1, -1)/\sqrt2 .
\]
The bands nearly always move together; the direction ``band 1 against band 2'' has almost no
variance. The condition number is $1.999/0.001 = 1999$.

\textbf{The inverse.} By \cref{eq:linalg:inv2},
\[
\Cov\inv = \frac{1}{0.001999}\begin{pmatrix} 1 & -0.999\\ -0.999 & 1\end{pmatrix}
= \begin{pmatrix} 500.25 & -499.75\\ -499.75 & 500.25\end{pmatrix}.
\]
Entries of 500, from a matrix whose entries were at most 1.

\textbf{Two pixels.} A pixel one unit up in both bands, $\vect{d} = (1, 1)$, has coordinate
$2/\sqrt2 = 1.414$ along $\vect{u}_1$ and 0 along $\vect{u}_2$:
$D^2 = 1.414^2/1.999 = 2/1.999 = 1.0005$. A pixel that differs by a tenth of a unit in opposite
directions, $\vect{d} = (0.1, -0.1)$, has coordinate 0 along $\vect{u}_1$ and $0.2/\sqrt2 = 0.1414$
along $\vect{u}_2$: $D^2 = 0.1414^2/0.001 = 0.02/0.001 = 20$. A difference ten times smaller in
Euclidean length scores twenty times higher.

That is RX working as designed, \emph{if} the eigenvalue 0.001 is right. Now suppose it was
estimated slightly wrong. If the true correlation is 0.999 but the estimate says 0.9995, the small
eigenvalue is $1 - 0.9995 = 0.0005$ and the same pixel scores $0.02/0.0005 = 40$; if the estimate
says 0.9999, it scores $0.02/0.0001 = 200$. A change in the fourth decimal place of one covariance
entry changed the score tenfold.
\end{worked}

Three things make small eigenvalues unreliable in practice, and all three bite hardest exactly where
RX looks hardest.

\paragraph{Sampling error.} A covariance estimated from a finite sample has errors in every entry.
The script draws 200 pixels from the $\rho = 0.999$ distribution of \cref{we:rx:nearsingular} five
times\tagC; the estimated small eigenvalue comes out as 0.00102, 0.00080, 0.00107, 0.00097 and
0.00093. Scores along that direction therefore vary by $\pm25\%$ from one sample to the next, while
the large eigenvalue's relative error is similar but harmless, because it divides a direction that
contributes little.

\paragraph{Rounding error.}\index{machine epsilon}\index{float32 and float64} A computer stores about 7 significant digits in \code{float32} (machine
epsilon $1.19\times10^{-7}$) and about 16 in \code{float64} ($2.22\times10^{-16}$). An eigenvalue
solver's absolute error is roughly epsilon times the largest eigenvalue, so any eigenvalue smaller
than about $\varepsilon\,\lambda_{\max}$ is rounding noise. A rule of thumb follows: inverting a
matrix with condition number $\kappa$ loses about $\log_{10}\kappa$ significant digits. For the
example, $\log_{10}1999 = 3.3$ digits. The script shows the extreme case: with $\rho = 1 - 10^{-8}$
the true small eigenvalue is $10^{-8}$; \code{float64} recovers $1.000000005\times10^{-8}$, but in
\code{float32} the number $1 - 10^{-8}$ rounds to exactly 1, the matrix becomes exactly singular, and
the small eigenvalue comes out as 0\tagC.

\paragraph{Dimension.} With 165 bands there are 165 eigenvalues, and neighbouring bands of a
reflectance spectrum are close to copies of each other, so many of the smallest eigenvalues sit at the
level of noise and estimation error rather than real variation. Each one is a direction like
$\vect{u}_2$ above, multiplied by its reciprocal. One such direction with an eigenvalue of
$10^{-10}$ and a pixel that differs along it by $10^{-3}$ contributes
$(10^{-3})^2/10^{-10} = 10^{-6}/10^{-10} = 10^{4}$ to the score, more than every meaningful direction
together (\cref{ch:residualmaha} works the same arithmetic). The ranking is then decided by which
pixels happen to have noise along the noisiest directions. That is what ``distances blow up to
1e11''\tagM{backend/allotrope/foundation_models/resolver.py} means: the numbers are not merely large,
they are meaningless.

A careful reader will object using \cref{sec:rx:chi2}: in exact arithmetic the in-sample scores
average exactly $C(N-1)/N$, about 165, so how can any be $10^{11}$? The identity holds only when the
covariance that is inverted is exactly the covariance of the pixels that are scored. Three things break
that link in practice, each through the small eigenvalues: rounding in the inverse (a direction whose
true eigenvalue is below $\varepsilon\,\lambda_{\max}$ gets an arbitrary tiny or zero estimate, as in the
\code{float32} example), statistics estimated on different pixels from those scored (a ring, a keep mask,
a sample), and values substituted into the cube (the detectors fill a pixel's missing bands with the band
mean, which puts it off the scene's correlation structure in exactly the low-variance directions). The
anomaly score is then dominated by the mismatch, not by the anomaly.

\begin{keypoint}
Inverting a covariance turns its smallest, least reliable eigenvalues into its largest weights; on
165 correlated bands those directions are noise, and RX ranks pixels by noise.
\end{keypoint}

\section{Regularisation}\label{sec:rx:regularisation}

The standard repair is to make the small eigenvalues less small before inverting.

\begin{definition}[Diagonal loading]\index{diagonal loading}\index{regularisation}
\emph{Diagonal loading} (also called ridge regularisation) replaces $\Cov$ by $\Cov + \epsilon\mat{I}$
for a small $\epsilon > 0$ before inverting.
\end{definition}

Its effect on the eigen-decomposition is exact and short. If $\Cov\vect{u}_k = \lambda_k\vect{u}_k$,
then
\[
(\Cov + \epsilon\mat{I})\vect{u}_k = \Cov\vect{u}_k + \epsilon\vect{u}_k = \lambda_k\vect{u}_k +
\epsilon\vect{u}_k = (\lambda_k + \epsilon)\vect{u}_k .
\]
The eigenvectors are unchanged and every eigenvalue moves up by $\epsilon$. Three consequences
follow.
\begin{enumerate}
  \item The matrix is invertible even when $\Cov$ is singular: the smallest eigenvalue is now at
    least $\epsilon$.
  \item The condition number falls to $(\lambda_{\max} + \epsilon)/(\lambda_{\min} + \epsilon)$.
  \item In the score, each direction contributes $(\vect{u}_k\T\vect{d})^2/(\lambda_k + \epsilon)$
    instead of $(\vect{u}_k\T\vect{d})^2/\lambda_k$. Directions with $\lambda_k \gg \epsilon$ are almost
    untouched; directions with $\lambda_k \ll \epsilon$ are capped at
    $(\vect{u}_k\T\vect{d})^2/\epsilon$.
\end{enumerate}

\begin{worked}[label={we:rx:loading}]{Loading the correlated pair}
For $\Cov = \begin{psmallmatrix} 1 & 0.999\\ 0.999 & 1\end{psmallmatrix}$ (eigenvalues 1.999 and
0.001) and the pixel $\vect{d} = (0.1, -0.1)$, whose only coordinate is 0.1414 along $\vect{u}_2$
(squared: 0.02):
\begin{center}
\begin{tabular}{cccc}
\toprule
$\epsilon$ & eigenvalues & condition number & $D^2 = 0.02/(\lambda_2 + \epsilon)$\\
\midrule
0         & 1.999, 0.001   & $1.999/0.001 = 1999$   & $0.02/0.001 = 20$\\
$10^{-4}$ & 1.9991, 0.0011 & $1.9991/0.0011 = 1817$ & $0.02/0.0011 = 18.2$\\
$10^{-3}$ & 2.000, 0.002   & $2.000/0.002 = 1000$   & $0.02/0.002 = 10$\\
$10^{-2}$ & 2.009, 0.011   & $2.009/0.011 = 183$    & $0.02/0.011 = 1.82$\\
\bottomrule
\end{tabular}
\end{center}
The large eigenvalue barely moves; the small one is rescued. The choice of $\epsilon$ decides how
much of the ``bands disagreeing'' direction you still believe: at $10^{-4}$ nearly all of it, at
$10^{-2}$ almost none.
\end{worked}

Two warnings come with diagonal loading.

\emph{The right $\epsilon$ depends on the units.} $\epsilon$ is added to variances, so it is in
squared data units. $10^{-4}$ is huge for reflectances whose variances are around $10^{-4}$, and
negligible for MNF components whose variances are 1 or more (\cref{ch:localrx}). A fixed default is
only meaningful for data on a known scale.

\emph{It keeps the solve from failing; it does not restore meaning.} With 165 near-degenerate
directions, a small $\epsilon$ leaves them nearly as amplified as before, and a large one throws away
the very directions RX is meant to use. The anomaly course puts it this way: the ridge ``keeps the
solve from failing but does not restore meaning to 165 near-degenerate dimensions''. The product's
fix (\cref{ch:localrx}) is to choose a small number of well-estimated directions first, and only
then to whiten.

\paragraph{Shrinkage.} A more general family, \emph{shrinkage}\index{shrinkage}, blends the
covariance with a multiple of the identity:
$(1 - \alpha)\Cov + \alpha\,\bar\lambda\,\mat{I}$, with $\bar\lambda = \tr\Cov/C$ the average
eigenvalue and $0 \le \alpha \le 1$. By the same argument the eigenvectors are unchanged and each
eigenvalue becomes $(1 - \alpha)\lambda_k + \alpha\bar\lambda$: pulled towards the average, small ones
up and large ones down. For the three-band $\Cov$ of \cref{we:rx:three} (average eigenvalue
$6/3 = 2$) with $\alpha = 0.1$, the eigenvalues $3.414, 2, 0.586$ become
$0.9\times3.414 + 0.2 = 3.273$, $0.9\times2 + 0.2 = 2$ and $0.9\times0.586 + 0.2 = 0.727$. Unlike
diagonal loading, shrinkage scales with the data automatically. No detector in this repository uses
it; it is here so that you recognise it in the literature.

\begin{sourcenote}[: what the code actually does to small eigenvalues]
Four different safeguards appear in the code, and only one of them is diagonal loading.
\begin{itemize}
  \item Global RX (\code{GlobalRXDetector}, via \code{spectral.rx}): none. It forms
    $\hat\Cov\inv$ with \code{np.linalg.inv}. It logs a warning when there are fewer than 10 pixels
    per band (\code{global_rx_detector.py}, lines 145--150), which addresses sample size, not
    conditioning.
  \item Local RX and MNF local RX: diagonal loading with $\epsilon = 10^{-4}$
    (\code{DEFAULT_REGULARIZATION = 1e-4}, \code{local_rx_detector.py} line 55 and
    \code{mnf_compression_lrx_detector.py} line 39; added in
    \code{app/detectors/_local_background.py} line 71), and a linear solve instead of an explicit
    inverse (line 78).
  \item MNF: a floor of $10^{-10}$ on the \emph{noise} eigenvalues
    (\code{mnf_compression_detector.py}, line 321), then truncation to 10 components.
  \item Residual Mahalanobis and PCA-RX: truncation to the top 30 principal directions and a floor of
    $10^{-12}$ on their eigenvalues (\code{residual_mahalanobis.py}, line 46).
\end{itemize}
The anomaly course says a ridge term ``is used, in the local variants''; that matches the code.
\end{sourcenote}

\paragraph{Solving instead of inverting.}\index{linear solve} The local detectors never form $\Cov\inv$. They solve the
linear system $(\Cov + \epsilon\mat{I})\vect{s} = \vect{d}$ for $\vect{s}$ with
\code{torch.linalg.solve}, then take $D^2 = \vect{d}\T\vect{s}$. The two are equal in exact arithmetic,
because $\vect{s} = \Cov\inv\vect{d}$; in floating point, solving is faster and loses fewer digits
than computing every entry of the inverse and then multiplying. The source comment says it plainly:
``Solving beats forming the inverse explicitly - better conditioned and faster.''

\section{The thermal case: RX in one band}\label{sec:rx:thermal}
\index{thermal RX}

With one band, $\Cov$ is the single number $\sigma^2$, its inverse is $1/\sigma^2$, and
\cref{eq:rx:rx} collapses to
\[
\text{RX}(x) = \frac{(x - \mu)^2}{\sigma^2} = z^2,
\]
the squared z-score you started the chapter with. For the thermal pixel, $3.42^2 = 11.7$, with a
$\chi^2_1$ tail probability of 0.000625\tagC, the same as the two-sided normal tail of 3.42. Single-band
covariance cannot be near-singular (unless the scene is perfectly uniform), which is why RX is safe on
thermal data and is shipped there as \code{thermal_grx}.

\begin{sourcenote}[: the thermal detector computes $z^2$ directly]
\code{ThermalGRXDetector.detect} (\code{app/detectors/thermal_grx_detector.py}, lines 95--104)
special-cases one band, because \code{spectral.rx} needs a covariance of at least $2\times2$: it
computes \code{((vals - mu) ** 2) / var} with \code{var = vals.var()}, numpy's divisor $N$, and
returns zeros if the variance is 0. Its mean and variance are those of the scene being scored, not the
shipped constants in \code{thermal_pixel_consts.json} (those normalise inputs for the thermal neural
model, Chakshu); the course's worked numbers use the shipped constants as a stand-in for a typical
scene.
\end{sourcenote}

\begin{confusions}
\confusion{Is RX machine learning?}{No. It is descriptive statistics computed from the scene in front
of it: a mean, a covariance, a quadratic form. No training, no checkpoint, no labels. Hand it a new
sensor and it works immediately.}
\confusion{Is a high RX score a large difference from the mean?}{Not necessarily. Pixel B of
\cref{we:rx:zfail} and pixel $(12,19,32)$ of \cref{we:rx:three} are smaller in Euclidean terms than
pixels that score lower. RX rewards differences in directions the background rarely takes.}
\confusion{Why is the Mahalanobis distance squared?}{Strictly, the Mahalanobis distance is the square
root of $\vect{d}\T\Cov\inv\vect{d}$. The code and this book work with the square, $D^2$, because that is
what appears in the Gaussian's exponent and what follows the chi-square distribution. The ranking is
the same either way.}
\confusion{Can I threshold RX at the chi-square value for 1\%?}{Only if the background really is one
Gaussian, which a real scene is not; you would flag far more than 1\% of pixels. The product thresholds
every score by percentile of its own map.}
\confusion{Why not regularise the 165-band covariance and keep plain RX?}{A ridge keeps the solve from
failing but cannot tell real low-variance directions from estimation noise; it either keeps the noise
(small $\epsilon$) or discards the real directions too (large $\epsilon$). MNF-RX and the top-30 PCA scores
choose well-estimated directions first.}
\confusion{Does a larger number of bands make RX more sensitive?}{It makes scores larger (the mean is
$C$) and the covariance harder to estimate (about $C^2/2$ numbers). More bands help only while each new
direction is estimated well; on 165 correlated bands they do not.}
\end{confusions}

\begin{checkyourself}
\item\label{ex:rx:1} A thermal pixel reads 61\,\textdegree C; use the scene mean 24.58 and standard
  deviation 13.57. Compute its z-score and its one-band RX score. Is it beyond the $\chi^2_1$ threshold
  for a 1\% false-alarm rate?
\item\label{ex:rx:2} For $\Cov = \begin{psmallmatrix} 5 & 2\\ 2 & 2\end{psmallmatrix}$ (the matrix of
  \cref{ex:linalg:3}), compute $\Cov\inv$ and the Mahalanobis distances of $\vect{d} = (1, -2)$ and
  $\vect{d} = (2, 1)$. Both have Euclidean length $\sqrt5$; explain the difference using the eigenvectors.
\item\label{ex:rx:3} With the three-band $\Cov$ and $\Cov\inv$ of \cref{we:rx:three}, score the deviation
  $\vect{d} = (0, 2, 0)$ (only the middle band is high). Compare with its per-band score.
\item\label{ex:rx:4} Two bands with variance 1 and correlation 0.99. Find the eigenvalues and the
  condition number, the Mahalanobis distance of $\vect{d} = (0.1, -0.1)$, and both again after diagonal
  loading with $\epsilon = 0.01$.
\item\label{ex:rx:5} An MNF-RX score map has $C = 10$ and is computed on a $1000\times1000$ scene. If the
  background were exactly Gaussian, how many pixels would exceed 23.21? Name two reasons a real scene
  exceeds it far more often.
\item\label{ex:rx:6} Show that for $C = 2$ the threshold with false-alarm rate $\alpha$ is $2\ln(1/\alpha)$,
  and compute it for $\alpha = 0.1\%$.
\end{checkyourself}

\nextchapter{ch:localrx}
